\section{Two pointers y sliding window}
\whenuse{Subarreglos/pares ordenados con una propiedad mon\'otona al mover los
extremos, pudiendo agregar y quitar elementos en $O(1)$.}

Si al avanzar el extremo derecho la validez cambia mon\'otonamente, cada
puntero avanza a lo sumo $n$ veces.

\complexity{$O(n)$ si agregar/quitar cuesta $O(1)$; memoria $O(1)$ o la del
estado mantenido.}

\begin{lstlisting}[style=cpp]
// Longitud maxima con suma <= limit; requiere a[i] >= 0.
int longest_window(const vector<ll> &a, ll limit) {
    ll sum = 0;
    int ans = 0;
    for (int l = 0, r = 0; r < (int)a.size(); ++r) {
        sum += a[r];
        while (l <= r && sum > limit) sum -= a[l++];
        ans = max(ans, r - l + 1);
    }
    return ans;
}

// Pares con suma target en un arreglo ordenado.
bool has_pair_sum(const vector<ll> &a, ll target) {
    int l = 0, r = (int)a.size() - 1;
    while (l < r) {
        i128 sum = (i128)a[l] + a[r];
        if (sum == target) return true;
        if (sum < target) ++l;
        else --r;
    }
    return false;
}

// Longitud maxima con suma <= limit; admite valores negativos.
int longest_window(const vector<ll>& a, ll limit) {
    vector<ll> maximum_prefix{0};
    vector<int> first_position{0};
    ll prefix = 0;
    int answer = 0;
    for (int r = 0; r < (int)a.size(); ++r) {
        prefix += a[r];
        // Buscamos el menor i con prefix[i] >= prefix[r+1] - limit.
        ll required = prefix - limit;
        auto it = lower_bound(
            maximum_prefix.begin(), maximum_prefix.end(), required
        );
        if (it != maximum_prefix.end()) {
            int k = it - maximum_prefix.begin();
            answer = max(answer, r + 1 - first_position[k]);
        }
        // Guardamos solamente nuevos maximos prefijos.
        if (prefix > maximum_prefix.back()) {
        maximum_prefix.push_back(prefix);
            first_position.push_back(r + 1);
        }
    }
    return answer;
}
\end{lstlisting}
